% 问题四：求解过程

\textbf{4 种方法求解结果对比如表~\ref{tab:q4_compare} 所示。}

\begin{table}[htbp]
\centering
\caption{问题 4 四种方法求解结果对比}
\label{tab:q4_compare}
\begin{tabular}{crrrrr}
\toprule
方法 & 撤销数 & A 撤销 & B 撤销 & C 撤销 & 调整数 \\
\midrule
贪心+间隔维度 & \textbf{44} & 0 & 1 & 43 & 73 \\
模拟退火+间隔维度 & \textbf{44} & 0 & 1 & 43 & 73 \\
贪心+局部搜索+间隔维度 & \textbf{44} & 0 & 1 & 43 & 73 \\
整数规划 ILP+间隔动作 & 46 & 1 & 10 & 35 & 84 \\
\bottomrule
\end{tabular}
\end{table}

\textbf{方案 1：贪心 + 优先级 + 间隔维度}

在问题 2 方案 1 的基础上，为 C 类装备新增间隔调整动作。动作优先级为：保留 → 频段平移（$\pm 1$ 到 $\pm 10$）→ 时间平移（$\pm 1$ 到 $\pm 5$）→ 间隔调整（仅 C 类，$\pm 1$ 到 $\pm 10$）→ 撤销。每档内按幅度从小到大枚举，选择最小幅度的可行动作。

结果：撤销 44 个（A:0, B:1, C:43），调整 73 个（其中 12 个为间隔调整），消解后冲突数为 0。

\textbf{方案 2：模拟退火 + 间隔维度}

以方案 1 的贪心解为初始解，采用以下策略：
\begin{itemize}
    \item 随机扰动：每次随机挑选 1\textasciitilde3 个计划，为其更换合法动作（包括间隔调整）；
    \item 多种子策略：使用多个不同随机种子各独立运行，取最优结果；
    \item 收尾消毒：退火结束后逐对检测残余冲突，撤销每对中优先级较低的计划。
\end{itemize}

结果：与方案 1 完全相同（撤销 44，调整 73）。说明贪心解已接近局部最优，模拟退火未能进一步改善。

\textbf{方案 3：贪心 + 局部搜索 + 间隔维度}

在方案 1 的贪心解基础上进行两阶段局部改进：一是\textbf{撤销恢复}，逐个尝试把被撤销的计划重新安置（允许使用间隔调整动作），若不引起新冲突则救回；二是\textbf{调整瘦身}，多轮扫描所有被调整的计划，尝试回退为保留或更小幅度的动作。由于贪心解已较优，两个阶段均无改进空间（救回 0 个、瘦身 0 处）。

结果：与方案 1 完全相同（撤销 44，调整 73）。

\textbf{方案 4：整数线性规划 ILP + 间隔动作}

将间隔调整作为新的动作类型纳入 ILP 模型，以撤销数最小化为目标函数，用 CBC 求解器求解。由于动作空间进一步扩大（新增 21 种间隔调整动作），求解难度增加。

结果：撤销 46 个（A:1, B:10, C:35），效果最差。ILP 效果差的原因在于动作空间过大导致松弛困难，求解器时间限制内无法找到高质量解。

\textbf{最优方案选择}

贪心+间隔维度、模拟退火+间隔维度、贪心+局部搜索+间隔维度三种方法的结果完全相同，撤销数均为 44 个，且 A 类撤销为 0，符合优先级要求。以贪心+间隔维度方案作为问题 4 的最终输出。

\textbf{与问题 2 对比}

\begin{table}[htbp]
\centering
\caption{问题 4 与问题 2 最优结果对比}
\label{tab:q4_vs_q2}
\begin{tabular}{crrr}
\toprule
指标 & 问题 2（SA） & 问题 4（贪心+间隔） & 改善 \\
\midrule
撤销总数 & 50 & \textbf{44} & -6 \\
A 撤销 & 0 & 0 & 0 \\
B 撤销 & 3 & 1 & -2 \\
C 撤销 & 47 & 43 & -4 \\
调整数 & 65 & 73 & +8 \\
总幅度 & 393 & 350 & -43 \\
\bottomrule
\end{tabular}
\end{table}

放宽 C 类装备的间隔约束后，撤销总数从 50 个减少到 44 个（减少 6 个），其中 B 类减少 2 个、C 类减少 4 个。调整数从 65 个增加到 73 个（增加 8 个），但总调整幅度从 393 降低到 350。这表明间隔调整动作为冲突消解提供了更多自由度，使得部分原本需要撤销的计划可以通过调整间隔保留下来。

消解后统计结果如表~\ref{tab:q4_result} 所示。

\begin{table}[htbp]
\centering
\caption{放宽间隔约束的时频冲突消解统计结果}
\label{tab:q4_result}
\begin{tabular}{crrr}
\toprule
类别 & 保留数量 & 调整数量 & 撤销数量 \\
\midrule
A 类 & 19 & 1 & 0 \\
B 类 & 7 & 32 & 1 \\
C 类 & 7 & 40 & 43 \\
\bottomrule
\end{tabular}
\end{table}
